% ============================================================
% Chapter III: Methodology
% ============================================================
\clearpage
\section{METHODOLOGY}

\subsection{System Architecture Overview}

The proposed framework consists of three interconnected layers: the simulation environment layer, the data acquisition layer, and the reconstruction and planning layer. Figure \ref{fig:system_architecture} illustrates the overall system architecture and data flow.

\begin{figure}[htbp]
\centering
% Placeholder for system architecture diagram
\fbox{\parbox{0.8\textwidth}{\centering
\vspace{2cm}
System Architecture Diagram\\
(UE5 + ASVSim Simulation Environment)\\
$\downarrow$\\
(Multi-Modal Data Collection Pipeline)\\
$\downarrow$\\
(3DGS Reconstruction + Path Planning)
\vspace{2cm}
}}
\caption{System architecture showing the three-layer framework integrating simulation environment, data collection, and reconstruction/planning modules.}
\label{fig:system_architecture}
\end{figure}

The simulation environment layer, built upon Unreal Engine 5 and ASVSim, provides realistic physics-based vessel dynamics and high-fidelity rendering. The data acquisition layer implements automated collection of synchronized multi-modal sensor data. The reconstruction and planning layer processes the collected data for 3D Gaussian Splatting training and supports path planning algorithms.

\subsection{Simulation Environment Configuration}

\subsubsection{ASVSim and Unreal Engine 5 Integration}

ASVSim \parencite{schroeder2020asvsim} is an autonomous surface vehicle simulator built on Unreal Engine 5, providing a Gymnasium-compatible reinforcement learning interface alongside sensor simulation capabilities. The simulator implements physics-based vessel dynamics using the Fossen model \parencite{fossen2011handbook}, which captures the hydrodynamic forces acting on vessels in calm water.

The simulation environment configuration is defined through a \texttt{settings.json} file that specifies:

\begin{itemize}
    \item Simulation mode and physics engine parameters
    \item Vehicle properties and initial conditions
    \item Sensor configurations (cameras, LiDAR)
    \item Rendering and capture settings
\end{itemize}

Key simulation parameters are listed in Table \ref{tab:simulation_params}.

\begin{table}[htbp]
\centering
\caption{ASVSim Simulation Configuration Parameters}
\label{tab:simulation_params}
\begin{tabular}{@{}lll@{}}
\toprule
\textbf{Parameter} & \textbf{Value} & \textbf{Description} \\
\midrule
SimMode & Vessel & Simulation mode \\
PhysicsEngine & VesselEngine & Physics engine type \\
ViewMode & SpringArmChase & Camera view mode \\
ClockSpeed & 1.0 & Real-time simulation \\
VehicleType & MilliAmpere & Vessel model \\
HydroDynamics & FossenCurrent & Hydrodynamics model \\
\bottomrule
\end{tabular}
\end{table}

\subsubsection{Sensor Configuration}

The sensor configuration defines the multi-modal perception capabilities of the simulated vessel. Two primary sensors are configured: a front-facing camera and a top-mounted LiDAR.

\textbf{Camera Configuration:} The front camera is positioned at the vessel bow, providing forward-looking imagery for ice detection and navigation. Camera parameters were determined through systematic optimization to balance field of view, resolution, and computational efficiency.

\begin{table}[htbp]
\centering
\caption{Front Camera Configuration}
\label{tab:camera_config}
\begin{tabular}{@{}lll@{}}
\toprule
\textbf{Parameter} & \textbf{Initial} & \textbf{Optimized} \\
\midrule
Resolution & 1280$\times$720 & 640$\times$480 \\
FOV & 90$^\circ$ & 70$^\circ$ \\
Position X & 2.5m & 3.5m \\
Position Z & -1.8m & -2.5m \\
Pitch & -10$^\circ$ & 0$^\circ$ \\
Lumen GI & Enabled & Disabled \\
\bottomrule
\end{tabular}
\end{table}

The resolution was reduced from 1280$\times$720 to 640$\times$480 based on performance analysis showing that:

\begin{itemize}
    \item Higher resolution (1280$\times$720) required approximately 55--60 seconds per frame due to Lumen global illumination overhead
    \item Reduced resolution (640$\times$480) achieved approximately 6.5 seconds per frame---a \textbf{9$\times$ speed improvement}
    \item The 640$\times$480 resolution maintains sufficient quality for 3DGS training while enabling practical data collection
\end{itemize}

The field of view was narrowed from 90$^\circ$ to 70$^\circ$ to reduce scene complexity, and the camera position was raised and moved forward to exclude vessel structure from the field of view.

\textbf{LiDAR Configuration:} The top-mounted LiDAR provides 360$^\circ$ horizontal scanning with 16 vertical channels, generating 8,192 points per scan at 10 Hz.

\begin{table}[htbp]
\centering
\caption{LiDAR Configuration}
\label{tab:lidar_config}
\begin{tabular}{@{}lll@{}}
\toprule
\textbf{Parameter} & \textbf{Value} & \textbf{Description} \\
\midrule
NumberOfLasers & 16 & Vertical channels \\
PointsPerSecond & 300,000 & Point generation rate \\
RotationsPerSecond & 10 & Scan frequency \\
Range & 100.0m & Maximum detection range \\
Position Z & -1.0m & Deck-mounted height \\
DataFrame & SensorLocalFrame & Coordinate reference \\
\bottomrule
\end{tabular}
\end{table}

\subsubsection{Rendering Optimization}

Unreal Engine 5's Lumen global illumination system, while producing photorealistic lighting, significantly impacts SceneCapture performance. Analysis revealed that each image request triggered complete Lumen GI passes, consuming 30--60 seconds per pass at 1280$\times$720 resolution.

The optimization strategy disabled Lumen and other post-processing effects in the \texttt{CameraDefaults.CaptureSettings}:

\begin{lstlisting}[language={}, caption={Optimized Capture Settings}]
{
  "ImageType": 0,
  "Width": 640,
  "Height": 480,
  "FOV_Degrees": 70,
  "LumenGIEnable": false,
  "LumenReflectionEnable": false,
  "MotionBlurAmount": 0,
  "ChromaticAberrationScale": 0.0,
  "AutoExposureBias": 0.0
}
\end{lstlisting}

These optimizations reduced the per-frame acquisition time from approximately 55 seconds to 6--7 seconds while maintaining sufficient image quality for 3D reconstruction.

\subsection{Multi-Modal Data Collection Pipeline}

\subsubsection{Collection Strategy}

The data collection pipeline implements an automated trajectory-based acquisition strategy designed to provide multi-view coverage suitable for 3D Gaussian Splatting training. The vessel follows predefined trajectories while capturing synchronized sensor data at regular intervals.

Three trajectory modes are implemented:

\begin{enumerate}
    \item \textbf{Circular trajectory:} The vessel maintains constant thrust and angle settings, resulting in a circular path that provides 360$^\circ$ coverage of the surrounding environment.

    \item \textbf{Linear trajectory:} The vessel travels in a straight line for initial scene traversal, followed by turning maneuvers.

    \item \textbf{Random walk:} Directional perturbations are applied at regular intervals to increase trajectory diversity.
\end{enumerate}

The circular trajectory is the primary mode, using parameters:

\begin{itemize}
    \item Thrust: 0.3 (30\% of maximum)
    \item Angle: 0.6 (slight right turn; 0.5 = straight)
    \item Movement interval: 0.3 seconds between frames
\end{itemize}

\subsubsection{Data Synchronization}

Ensuring temporal synchronization across sensor modalities is critical for multi-modal fusion. The collection pipeline follows this sequence:

\begin{enumerate}
    \item Apply vessel control commands
    \item Wait for physics stabilization (0.3s)
    \item Capture RGB image and depth map (simultaneous \texttt{simGetImages} call)
    \item Query vessel state for pose information
    \item Acquire LiDAR scan
    \item Save all data with shared timestamp
\end{enumerate}

Notably, the \texttt{simPause} mechanism was disabled due to compatibility issues with CosysAirSim v3.0.1 that caused scene freezing. Instead, short movement intervals (0.3s) ensure minimal pose variation between sensor readings.

\subsubsection{Image Type Handling}

CosysAirSim uses different ImageType enumerations than standard AirSim, requiring careful configuration:

\begin{table}[htbp]
\centering
\caption{CosysAirSim ImageType Enumeration}
\label{tab:imagetype}
\begin{tabular}{@{}lll@{}}
\toprule
\textbf{ImageType} & \textbf{CosysAirSim} & \textbf{Standard AirSim} \\
\midrule
Scene (RGB) & 0 & 0 \\
DepthPlanar & \textbf{1} & \textbf{2} \\
DepthVis & 3 & 1 \\
Segmentation & 5 & 5 \\
\bottomrule
\end{tabular}
\end{table}

The critical difference is \texttt{DepthPlanar}, which uses value 1 in CosysAirSim rather than 2. This necessitated explicit configuration of all ImageType variants in \texttt{CameraDefaults.CaptureSettings}.

\subsubsection{Data Format Conversion}

Raw sensor data requires conversion to standardized formats:

\textbf{RGB Images:} CosysAirSim returns RGB-ordered data, which must be converted to BGR for OpenCV compatibility:

\begin{lstlisting}[language=Python, caption={RGB to BGR Conversion}]
def decode_image_safe(response, name=""):
    raw = np.frombuffer(response.image_data_uint8, dtype=np.uint8)
    h, w = response.height, response.width
    total = len(raw)

    if total == h * w * 4:  # RGBA
        img = raw.reshape(h, w, 4)[:, :, :3]
        img = cv2.cvtColor(img, cv2.COLOR_RGB2BGR)
    elif total == h * w * 3:  # RGB
        img = raw.reshape(h, w, 3)
        img = cv2.cvtColor(img, cv2.COLOR_RGB2BGR)
    return img
\end{lstlisting}

\textbf{Depth Maps:} Depth values are returned as float32 arrays in meters. Invalid depths (sky, distant objects) are represented as infinity.

\textbf{LiDAR Points:} Point clouds are stored in SensorLocalFrame as N$\times$3 numpy arrays.

\textbf{Pose Information:} Vessel pose includes position (X, Y, Z) and orientation (quaternion: w, x, y, z) from the ASVSim world coordinate system.

\subsection{Output Data Structure}

Collected datasets are organized in timestamped directories with the following structure:

\begin{lstlisting}[language=bash, caption={Dataset Directory Structure}]
dataset/YYYY_MM_DD_HH_MM_SS/
|-- rgb/              # RGB images (PNG, BGR format)
|   |-- 0000.png
|   |-- 0001.png
|   `-- ...
|-- depth/            # Depth maps (NPY, float32, meters)
|   |-- 0000.npy
|   |-- 0001.npy
|   `-- ...
|-- lidar/            # LiDAR point clouds (JSON)
|   |-- 0000.json
|   |-- 0001.json
|   `-- ...
|-- meta/             # Metadata and configuration
|   |-- collection_config.json
|   `-- collection_report.json
|-- colmap/           # COLMAP format for 3DGS
|   `-- transforms.json
`-- poses.json        # Camera poses (ground truth)
\end{lstlisting}

\subsubsection{Pose File Format}

The \texttt{poses.json} file contains per-frame pose and camera intrinsics:

\begin{lstlisting}[language={}, caption={Pose JSON Structure}]
{
  "frame_id": 0,
  "position": {"x": 1.2, "y": 0.3, "z": 0.01},
  "orientation": {"w": 0.99, "x": 0.0, "y": 0.0, "z": 0.1},
  "speed": 0.45,
  "camera_intrinsics": {
    "fx": 466.0, "fy": 466.0,
    "cx": 320.0, "cy": 240.0,
    "width": 640, "height": 480,
    "fov_deg": 70.0
  },
  "timestamp": "2026-03-12T10:30:00"
}
\end{lstlisting}

Camera intrinsics are computed offline using the pinhole camera model:

\begin{equation}
f_x = f_y = \frac{W}{2 \tan(\text{FOV}/2)}
\label{eq:intrinsics}
\end{equation}

\begin{equation}
c_x = W/2, \quad c_y = H/2
\label{eq:principal_point}
\end{equation}

where $W$ and $H$ are image width and height, and FOV is the field of view in radians.

\subsection{COLMAP Format Conversion}

For 3D Gaussian Splatting training, poses are converted to the transforms.json format used by NeRF and 3DGS implementations:

\begin{lstlisting}[language={}, caption={COLMAP Transforms Format}]
{
  "camera_model": "OPENCV",
  "fl_x": 466.0, "fl_y": 466.0,
  "cx": 320.0, "cy": 240.0,
  "w": 640, "h": 480,
  "frames": [
    {
      "file_path": "rgb/0000.png",
      "transform_matrix": [[...], [...], [...], [...]]
    }
  ]
}
\end{lstlisting}

The transformation matrix converts from camera coordinates to world coordinates. For simulation data with known ground truth, this enables direct training without Structure-from-Motion preprocessing.

\subsection{Error Handling and Validation}

The data collection pipeline implements robust error handling:

\textbf{Retry Mechanism:} Failed frame acquisitions are automatically retried up to 3 times with 1-second delays.

\textbf{Frame Validation:} Each frame undergoes validation checks:

\begin{itemize}
    \item Image dimensions match expected resolution
    \item Depth map contains finite values within expected range
    \item LiDAR point cloud contains expected number of points (8,192)
    \item Pose data is finite and within scene bounds
\end{itemize}

\textbf{Dataset Validation:} Post-collection validation verifies:

\begin{itemize}
    \item Consistent file counts across modalities
    \item Sequential frame numbering
    \item Pose continuity (no large jumps between frames)
    \item Image quality (proper color channels, no corruption)
\end{itemize}

\subsection{Implementation Details}

The data collection system is implemented in Python using the CosysAirSim SDK (v3.0.1). Key implementation characteristics:

\begin{itemize}
    \item \textbf{No multi-threading:} Avoids msgpack-rpc IOLoop errors
    \item \textbf{No simPause:} Prevents scene freezing in CosysAirSim v3.0.1
    \item \textbf{Batch API calls:} Multiple image types requested in single \texttt{simGetImages} call
    \item \textbf{Graceful shutdown:} Proper vessel control release on interruption
    \item \textbf{Progress tracking:} Real-time display of frame count, position, and timing
\end{itemize}

The complete source code for the data collection pipeline is provided in Appendix A.
